% !TeX root = ../prompt-engineering-guide-for-beginners.tex

%%%%%%%%%%%%%%%%%%%%%%%%%%%%%%%%%%
%% 第十章：多模态提示
%%%%%%%%%%%%%%%%%%%%%%%%%%%%%%%%%%
\chapter{多模态交互：不只是打字聊天}
\label{ch:multimodal}

到目前为止，我们聊的都是文字文本层面的对话。但现在的顶级大语言模型进化极快，它们早已经不仅仅局限于会「读」和「写」了：它们还能看懂图片、透视复杂文档、听懂语音，甚至挥毫泼墨帮你画画。
这一章将带你领略这些被统称为「多模态（Multimodal）」的能力，以及如何使用专属的提示词点亮它们。

\section{看图说话：图片理解与分析}

越来越多的 AI 模型支持直接上传图片进行分析。你可以让 AI 做这些事：

\begin{itemize}
  \item \textbf{描述图片内容}：上传一张照片，让 AI 告诉你图里有什么
  \item \textbf{提取图片中的文字}（OCR）：拍一张名片、截一张表格，让 AI 把里面的文字提取出来
  \item \textbf{分析图表数据}：上传柱状图、折线图、饼图等，让 AI 帮你解读数据趋势
  \item \textbf{辅助设计}：上传一张界面截图，让 AI 分析布局并提出改进建议
\end{itemize}

写图片分析提示词的技巧：

\begin{promptbox}
请看这张图片并完成以下任务：
\begin{enumerate}
  \item 描述图片中的主要内容
  \item 提取图片中所有可见的文字信息
  \item 如果是图表，请分析数据并给出 3 个关键发现
\end{enumerate}
\end{promptbox}

关键要点是：\textbf{千万不要偷懒只说一句「看看这张图」}。就像文字提示词一样，你需要明确界定你希望 AI 扮演什么角色、从图片中抓取什么信息、基于图片做什么具体分析。你的指令越有指向性，它的“视力”与“见解”就越敏锐。

\section{喂文档：让 AI 闪电阅读你的文件}

许多优秀的 AI 工具现在已原生支持直接上传外部文件（包括 PDF、Word、Excel、PPT 乃至长篇代码文件等），并可直接让 AI 阅读、交叉比对和总结回答其中的问题。这比你自己吭哧吭哧开文档去复制粘贴要优雅无数倍。

常见的文档处理神仙级应用场景：

\begin{itemize}
  \item \textbf{快速总结}：「请总结这份 PDF 的核心观点，不超过 300 字。」
  \item \textbf{文档问答}：「根据这份合同内容，乙方的违约责任具体是什么？」
  \item \textbf{多文档对比}：「请对比这两份报价单，找出价格差异超过 10\% 的项目。」
  \item \textbf{数据处理}：上传 Excel 文件，让 AI 帮你做统计、生成图表描述、找出异常值
\end{itemize}

使用文档功能时的提示词技巧：

\begin{promptbox}
我上传了一份公司年度财务报告（PDF）。请完成以下任务：
\begin{enumerate}
  \item 提取报告中的关键财务指标（营收、净利润、毛利率）
  \item 与上一年度进行对比（如果报告中有提及）
  \item 用 3 到 5 句话总结公司的财务健康状况
  \item 如果你发现报告中有值得注意的风险点，请单独列出
\end{enumerate}
\end{promptbox}

\begin{tipbox}
受限于算力和接口限制，许多工具对单次上传的文件体量（MB 数和总页数）设有硬性天花板。遇到动辄大几百页的超级年报，为了保证它检索的准确度，强烈建议拆分成多个章节文档按需分别投喂。
\end{tipbox}

\section{语音交互：用嘴代替键盘}

很多 AI 工具已经支持语音输入和语音输出。你可以直接用说话的方式和 AI 对话，不用打字。

语音交互最大的价值在于「解放双手」。想象一下这些场景：你正在开车上班的路上，突然想到一个文案的灵感，但没法打字；或者你在做饭时想让 AI 帮你整理一下明天的会议议程；又或者你想边散步边和 AI 练习英语口语。这些场景下，语音交互都比打字方便得多。

语音输入能为你带来「口述式写作」的奇特爽感。你在手写文案时可能会长时间卡壳，但嘴巴说起来往往滔滔不绝。你大可直接对着手机 App 版的模型把脑内繁杂琐碎的想法全都宣泄出来，最后收尾补上一句「请把刚才这些散乱的想法帮我整理成一篇结构严谨的述职大纲」。这对不擅长打底稿的人来说，无异于降维打击。

使用语音交互的心法小窍门：

\begin{itemize}
  \item 语意尽量连贯完整。不要说半句话就长久停顿，AI 极可能会将这半句话立刻当成完整指令去执行，导致结果错位。
  \item 在人名与数字等极精密的专有名词上放慢咬字。例如读错「三十五十万」和「三千五十万」将会产生巨大的业务灾难。
  \item 万一嘴瓢说错了，无需终止重来。你可以直接像和真人聊天一样补充纠正：「哎呀不对，我上面说的不是 X 而是 Y，请按 Y 作为标准」。AI 的语义容错理解能力足以完美剥离出正确的意图。
  \item 有节奏的留白。在表达完一个小节要点后维持一秒短暂的空隙，能大幅增强内置语音识别系统抓取断句标点符号的准确性。
\end{itemize}

\section{AI 生图：用文字画画}

现在有很多 AI 可以根据你的文字描述生成图片（叫做「文生图」），比如 DALL·E、Midjourney、Stable Diffusion 等。

写图片生成提示词和写文字提示词有一些不同，重点在于对视觉元素的描述。

一个好的生图提示词通常包含以下要素。你可以把它想象成给画家写的一张「创作需求单」，写得越具体，画家就越知道你想要什么：

\begin{itemize}
  \item \textbf{主体}：你想画什么？（一只猫、一座城市、一碗面）这是整个描述的基础，必不可少。
  \item \textbf{风格}：什么艺术风格？（水彩、油画、赛博朋克、吉卜力动画风格 、极简线条）这决定了图片的整体视觉感觉。
  \item \textbf{构图}：从什么角度看？（俯视、特写、全景、侧面）同一个对象从不同角度拍摄，给人的感觉完全不同。
  \item \textbf{光线和氛围}：什么样的光线和情绪？（暖色调、日落光线、雨天、温馨）光线和氛围往往是一张图从「还行」到「惊艳」的关键差别。
  \item \textbf{细节补充}：其他重要的视觉细节（背景模糊、高清、4K 等）
\end{itemize}

\textbf{示例：}

\begin{promptbox}
\textbf{简单版}：一只橘猫坐在窗台上

\textbf{详细版}：一只毛茸茸的橘猫坐在旧木窗台上，窗外是下雨的城市街景，画面风格为水彩插画，暖色调，柔和的自然光从窗户照进来，给人温馨治愈的感觉
\end{promptbox}

\begin{tipbox}
描述要素写得越详尽、越可视化，最终生成的图就越接近你的脑内设想。但请务必避免塞入太多「彼此互斥」的情境设定（例如要求「背景是正午的烈日」同时又要「充满昏暗阴郁的深夜惊悚感」），这会导致模型出现逻辑错乱甚至画出“克苏鲁”风格的缝合怪。
\end{tipbox}

\section{多模态提示词的特别红线}

在使用多模态功能时，有几点需要特别注意：

\begin{itemize}
  \item \textbf{模型能力差异}：不是所有 AI 都支持图片理解或文件上传。使用前确认你的工具支持哪些模态。
  \item \textbf{文件大小限制}：大多数工具对上传文件的大小有限制（通常几十 MB 以内），超大文件可能需要压缩或分割。
  \item \textbf{隐私保护}：上传图片和文档时要特别注意隐私。不要上传包含敏感个人信息、公司机密或他人隐私的内容。
  \item \textbf{图片中的文字}：AI 识别图片中文字的能力在不断提升，但仍然可能出错，特别是手写体、模糊图片或排版复杂的场景。重要信息建议人工核实。
  \item \textbf{生图的局限}：AI 生成的图片在细节上可能有瑕疵（比如手指数量不对、文字变形等），如果用于正式场合，需要仔细检查。
\end{itemize}

多模态能力是 AI 发展最快的方向之一。现在还做不到的事情，过几个月可能就能做了。保持对新功能的关注，适时尝试，你会发现越来越多意想不到的用法。
